\begin{thebibliography}{24}
\bibitem[1]{rc} Nantel Bergeron and Sara Billey, RC-graphs and Schubert polynomials, Exp. Math. 2 (1993), no. 4, 257–269.
\bibitem[2]{reviewedfansub} Nantel Bergeron, Cesar Ceballos, and Jean-Philippe Labbé, Fan realizations of type A subword complexes and multi-associahedra of rank 3, Discrete Comput. Geom. 54 (2015), no. 1, 195–231.
\bibitem[3]{cesar} Cesar Ceballos, On associahedra and related topics, Ph.D. thesis, Freie Universität Berlin (Germany), 2012.
\bibitem[4]{subwordcluster} Cesar Ceballos, Jean-Philippe Labbé, and Christian Stump, Subword complexes, cluster complexes, and generalized multi-associahedra, J. Algebr. Comb. 39 (2014), no. 1, 17–51.
\bibitem[5]{escobar} Laura Escobar, Brick manifolds and toric varieties of brick polytopes, Electron. J. Comb. 23 (2016), no. 2, article ID P2.25 (18 pages).
\bibitem[6]{fom-kir} Sergey Fomin and Anatol N. Kirillov, Grothendieck polynomials and the Yang-Baxter equation, in Formal power series and algebraic combinatorics, DIMACS, 1994, pp. 183–189.
\bibitem[7]{roothistory} Israel M. Gelfand, Mark I. Graev, and Alexander Postnikov, Combinatorics of hypergeometric functions associated with positive roots, in The Arnold-Gelfand mathematical seminars, Birkhäuser, 1997, pp. 205–221.
\bibitem[8]{darij} Darij Grinberg, t-unique reductions for Mészáros’s subdivision algebra, https://arxiv.org/ abs/1704.00839, 2017.
\bibitem[9]{k2014} Anatol N. Kirillov, On some quadratic algebras, Dunkl elements, Schubert, Grothendieck, Tutte and reduced polynomials, RIMS preprint, 2014.
\bibitem[10]{subword} Allen Knutson and Ezra Miller, Subword complexes in Coxeter groups, Adv. Math. 184 (2004), no. 1, 161–176.
\bibitem[11]{annals} Allen Knutson and Ezra Miller, Gröbner geometry of Schubert polynomials, Ann. Math. 161 (2005), no. 3, 1245–1318.
\bibitem[12]{root1} Karola Mészáros, Root polytopes, triangulations, and the subdivision algebra. I, Trans. Am. Math. Soc. 363 (2011), no. 8, 4359–4382.
\bibitem[13]{root2} Karola Mészáros, Root polytopes, triangulations, and the subdivision algebra, II, Trans. Am. Math. Soc. 363 (2011), no. 11, 6111–6141.
\bibitem[14]{h-poly1} Karola Mészáros, h-polynomials via reduced forms, Electron. J. Comb. 22 (2015), no. 4, article ID P4.18 (17 pages).
\bibitem[15]{prod} Karola Mészáros, Product formulas for volumes of flow polytopes, Proc. Am. Math. Soc. 143 (2015), no. 3, 937–954.
\bibitem[16]{h-poly2} Karola Mészáros, h-polynomials of reduction trees, SIAM J. Discrete Math. 30 (2016), no. 2, 736–762.
\bibitem[17]{groth} Karola Mészáros, Pipe dream complexes and triangulations of root polytopes belong together, SIAM J. Discrete Math. 30 (2016), no. 1, 100–111.
\bibitem[18]{mm} Karola Mészáros and Alejandro H. Morales, Flow polytopes of signed graphs and the Kostant partition function, Int. Math. Res. Not. (2015), no. 3, 830–871.
\bibitem[19]{groth1} Karola Mészáros and Avery St. Dizier, From generalized permutahedra to Grothendieck polynomials via flow polytopes, https://arxiv.org/abs/1705.02418, 2017.
\bibitem[20]{assoc} Vincent Pilaud and Michel Pocchiola, Multitriangulations, pseudotriangulations and primitive sorting networks, Discrete Comput. Geom. 48 (2012), no. 1, 142–191.
\bibitem[21]{abrick} Vincent Pilaud and Francisco Santos, The brick polytope of a sorting network, Eur. J. Comb. 33 (2012), no. 4, 632–662.
\bibitem[22]{SerranoStump} Luis Serrano and Christian Stump, Maximal fillings of moon polyominoes, simplicial complexes, and Schubert polynomials, Electron. J. Comb. 19 (2012), no. 1, article ID P16 (18 pages).
\bibitem[23]{Stcom} Richard P. Stanley, Combinatorics and commutative algebra, Progress in Mathematics, vol. 41, Birkhäuser, 1996.
\bibitem[24]{stump} Christian Stump, A new perspective on k-triangulations, J. Comb. Theory, Ser. A 118 (2011), no. 6, 1794–1800.
\end{thebibliography}
